% ============================================================================
% INTRODUCTION DU CHAPITRE 1 — CONTEXTE, \'ETAT DE L'ART ET PROBL\'eMATIQUE
% source : partie1/introduction_partie1.tex
% Version : 2.2 — alignée sur la formalisation Lean4 révisée
% ============================================================================

% ----------------------------------------------------------------------------
\section{INTRODUCTION}
\label{sec:introduction_contexte}
% ----------------------------------------------------------------------------

\begin{sloppypar}
En quatre d\'ecennies, la cryptographie prouvable a produit des garanties
d'une grande solidit\'e, qu'il s'agisse de la confidentialit\'e, mesur\'ee par
l'avantage d'un distingueur polynomial, ou de la fiabilit\'e, \'etablie par
r\'eduction \`a un probl\`eme computationnellement difficile. Ces r\'esultats,
exacts sur le plan math\'ematique, restent pourtant souvent inexploitables
lorsqu'une institution doit s'en saisir.

Un juge ne raisonne pas comme un distingueur polynomial, pas plus qu'un
auditeur financier ne v\'erifie un circuit arithm\'etique. Lorsque la preuve
cryptographique d'une transaction doit convaincre un r\'egulateur, ou qu'une
attestation zero-knowledge est vers\'ee au dossier d'une instruction p\'enale,
la question d\'ecisive devient celle de l'interpr\'etation de cette preuve
par un acteur institutionnel cognitivement born\'e. Cette th\`ese part de ce d\'eplacement de
perspective, dont le pr\'esent chapitre pose les fondements.

Trois syst\`emes cryptographiques d\'eploy\'es servent de construits
illustratifs, \`a savoir Helios pour le vote \'electronique, les zk-SNARKs
pour la conformit\'e financi\`ere et ECDSA pour les signatures
contractuelles. Leurs profils estim\'es laissent appara\^itre un motif que
l'\'etat de l'art interdisciplinaire permettra ensuite de nommer puis de
borner. La probl\'ematique est alors formul\'ee en six objectifs de
recherche qui structurent la progression canonique
$\mathbf{F}\to\mathbf{B}\to\mathbf{E}\to\mathbf{S}\to\mathbf{C}\to\mathbf{X}$
(Formaliser, Borner, \'Etendre, Saturer, Certifier, eXp\'erimenter) de
cette th\`ese. Apr\`es avoir situ\'e la recherche dans son cadre
institutionnel et professionnel (section~\ref{subsec:cadre_recherche}), nous
posons les fondements de la th\'eorie de l'information classique
(section~\ref{subsec:fondements_info}) avant de les \'etendre au cadre
quantique (section~\ref{subsec:fondements_info_quantique}). Les primitives
cryptographiques d'obscurcissement sont pr\'esent\'ees \`a la
section~\ref{subsec:primitives_crypto}, suivies d'une revue des trilemmes
d'impossibilit\'e fondateurs (section~\ref{subsec:trilemmes}). La
section~\ref{subsec:paradoxe_authenticite} formalise ensuite le paradoxe de
l'authenticit\'e invisible, dont la section~\ref{subsec:objectifs_recherche}
d\'erive les six objectifs de recherche et leurs contributions.
\end{sloppypar}